\chapter*{Introduction generale}
\addcontentsline{toc}{chapter}{Introduction generale}

Le contexte de ce projet est marque par la multiplication des cyberattaques visant les acces privilegies, alors que les modeles perimetriques traditionnels ne suffisent plus a proteger des environnements hybrides et distribues. Dans ce cadre, le principe Zero Trust, formalise notamment par le NIST SP 800-207, impose de ne jamais faire confiance par defaut et de verifier chaque demande d'acces.

La problematique adressee dans ce rapport est la suivante: comment supprimer les acces permanents tout en preservant la capacite operationnelle des equipes techniques? La reponse proposee repose sur une gouvernance JIT fondee sur l'authentification forte, l'approbation explicite, la limitation temporelle des privileges et la tracabilite systematique des actions.

La solution integre Authentik pour l'IAM, une application Flask pour la gouvernance des demandes, Tailscale ACL pour l'application des droits reseau, un moteur Python pour l'activation et la revocation temporaires, et Wazuh pour l'observabilite et l'audit.

Le rapport est structure en six chapitres. Le premier presente le contexte et les risques lies aux acces privilegies. Le deuxieme positionne la solution par rapport aux approches existantes. Le troisieme detaille la conception de l'architecture Zero Trust. Le quatrieme decrit l'implementation. Le cinquieme presente la validation et les tests. Le sixieme discute les limites et les travaux futurs, avant la conclusion generale et les annexes.

\clearpage